\documentclass[10pt]{article}
\usepackage[T1]{fontenc}
\usepackage[utf8]{inputenc}
\usepackage[french]{babel}
\usepackage{amsmath,amssymb,mathrsfs,mathpazo}
\usepackage{geometry}
\geometry{a4paper,margin=2cm}
\usepackage{graphicx,array,tabularx,tabularray,colortbl,multirow,multicol,booktabs,enumitem,pifont,ulem,fancybox,diagbox,cancel}
\usepackage{pstricks,pstricks-add,pst-plot,pst-tree,pst-node,pst-text,pst-func,pst-math,pst-bspline,pst-3dplot,pst-eucl,pst-solides3d}
\let\pscalculate\undefined
\usepackage{pst-bezier}
\usepackage{tikz,pgfplots,tkz-tab}
\usetikzlibrary{arrows,arrows.meta,patterns,calc,positioning,trees,decorations.markings,shapes.geometric,shapes.misc,shapes.arrows,intersections,decorations.pathmorphing,decorations.pathreplacing,backgrounds,shadows,babel}
\pgfplotsset{compat=1.18}
\usepgfplotslibrary{fillbetween}
\usepackage[np]{numprint}
\usepackage{listings,makecell,pas-tableur,twoopt}
\let\barre\undefined
\usepackage[pstricks]{tablvar}
\definecolor{darkgreen}{RGB}{0,100,0}
\providecommand{\red}{\color{red}}
\providecommand{\blue}{\color{blue}}
\providecommand{\green}{\color{green}}
\providecommand{\black}{\color{black}}
\providecommand{\white}{\color{white}}
\providecommand{\gray}{\color{gray}}
\providecommand{\cyan}{\color{cyan}}
\providecommand{\magenta}{\color{magenta}}
\providecommand{\yellow}{\color{yellow}}
\providecommand{\violet}{\color{violet}}
\providecommand{\orange}{\color{orange}}
\providecommand{\euro}{\texteuro}
\providecommand{\EUR}{\texteuro}
\providecommand{\ieme}{\textsuperscript{e}}
\providecommand{\degres}{\ensuremath{^\circ}}
\providecommand{\degree}{\ensuremath{^\circ}}
\providecommand{\vv}[1]{\overrightarrow{#1}}
\providecommand{\Coord}[1]{#1}
\providecommand{\decosix}{}
\providecommand{\starredbullet}{\ensuremath{\bullet}}
\pagestyle{empty}
\setlength{\parindent}{0pt}
\usepackage{adjustbox}\DeclareUnicodeCharacter{2236}{:}\DeclareUnicodeCharacter{2190}{\ensuremath{\leftarrow}}
\providecommand{\R}{\mathbb{R}}
\providecommand{\N}{\mathbb{N}}
\providecommand{\Z}{\mathbb{Z}}
\providecommand{\Q}{\mathbb{Q}}
\providecommand{\D}{\mathbb{D}}
\providecommand{\C}{\mathbb{C}}
\providecommand{\e}{{\mathrm{e}}}
\providecommand{\vect}[1]{\overrightarrow{#1}}
\providecommand{\vectt}[1]{\overrightarrow{\mathrm{#1}}}
\providecommand{\vvt}[1]{\overrightarrow{\mathrm{#1}}}
\providecommand{\np}[2][]{{#2~#1}}
\providecommand{\vv}[1]{{\overrightarrow{#1}}}
\renewcommand{\d}{{\mathrm{d}}}
\providecommand{\degres}{{^\circ}}
\providecommand{\degre}{{^\circ}}
\providecommand{\blue}{\color{blue}}
\providecommand{\red}{\color{red}}
\providecommand{\enskip}{\quad}
\providecommand{\llbracket}{\lbrack\!\lbrack}
\providecommand{\rrbracket}{\rbrack\!\rbrack}
\providecommand{\euro}{{\text{€}}}
\providecommand{\no}{{\text{n°}}}
\providecommand{\interval}[2]{{[#1;#2]}}
\providecommand{\intervFO}[2]{{[#1;#2[}}
\providecommand{\dotfill}{\ldots}

\providecommand{\ds}{\displaystyle}
\providecommand{\textding}[1]{#1}
\providecommand{\Oij}{\(\mathrm{O};\vec{\imath},\vec{\jmath}\)}
\providecommand{\Oijk}{\(\mathrm{O};\vec{\imath},\vec{\jmath},\vec{k}\)}
\providecommand{\Ouv}{\(\mathrm{O};\vec{u},\vec{v}\)}
\providecommand{\intervFF}[2]{\ensuremath{[#1;#2]}}
\renewcommand{\euro}{\ifmmode\text{\texteuro}\else\texteuro\fi}
\DeclareUnicodeCharacter{00B0}{\ifmmode^{\circ}\else\textdegree\fi}
\DeclareUnicodeCharacter{20AC}{\ifmmode\text{\texteuro}\else\texteuro\fi}
% Énoncés : Éducation nationale. Transcription et corrigés : APMEP, Denis Vergès et contributeurs. Découpage et classement : CoopMaths, Sébastien Lozano, Éric Elter et contributeurs. Adaptation de la présentation : Maths Explorer.
% CC BY-SA : https://coopmaths.fr/www/about/
% https://coopmaths.fr/alea/static/bac/2021/tex/bac_2021_01_sujet0_2.tex
\begin{document}
\noindent\textit{APMEP / CoopMaths. CC BY-SA. Présentation adaptée par Maths Explorer.}\par\medskip
\medskip

On considère le cube ABCDEFGH de côté 1, le milieu I de [EF] et J le symétrique de E par rapport à F{}.
\begin{center}

\psset{xunit=4cm,yunit=4cm}
\begin{pspicture}(-0.1,-0.1)(2,1.4)

\psframe(0,0)(1,1)
\psline(1,0)(1.35,0.35)
\psline[linestyle=dashed](0,0)(0.35,0.35)
\psline[linestyle=dashed](0.35,0.35)(1.35,0.35)
\psline[linestyle=dashed](0.35,0.35)(0.35,1.35)
\psline(1.35,0.35)(1.35,1.35)
\psline(0,1)(0.35,1.35)
\psline(0.35,1.35)(1.35,1.35)
\psline(1,1)(1.35,1.35)
\uput[dl](0,0){A}\uput[dr](1,0){B}\uput[r](1.35,0.35){C}
\uput[l](0.35,0.35){D}\uput[l](0,1){E}\uput[ul](1.05,1){F}
\uput[ur](1.35,1.35){G}\uput[ul](0.35,1.35){H}
\uput[u](0.5,1){I}

\psline[linewidth=0.8pt,linecolor=lightgray] (1,1)(2,1)\rput(2.05,1){J}
\psdots(0,0)(1,1)(0,1)(1,0)(1.35,0.35)(0.35,0.35)(0.5,1)(2,1)(1.35,1.35)(0.35,1.35)
\end{pspicture}

\end{center}

\medskip

Dans tout l'exercice, l'espace est rapporté au repère orthonormé $\left(\text{A}~;~\vv{\text{AB}},~\vv{\text{AD}},~\vv{\text{AE}}\right)$.

\medskip

\begin{enumerate}
\item 
	\begin{enumerate}
		\item Par lecture graphique, donner les coordonnées des points I et J.
		\item En déduire les coordonnées des vecteurs $\vv{\text{DJ}}$ , $\vv{\text{BI}}$ et $\vv{\text{BG}}$.
		\item Montrer que $\vv{\text{DJ}}$ est un vecteur normal au plan (BGI).
		\item Montrer qu’une équation cartésienne du plan (BGI) est $2x - y + z - 2 = 0$.
	\end{enumerate}
\item On note $d$ la droite passant par F et orthogonale au plan (BGI).
	\begin{enumerate}
		\item Déterminer une représentation paramétrique de la droite $d$.
		\item On considère le point L de coordonnées $\left(\frac{2}{3}~;~\frac{1}{6}~;~\frac{5}{6}\right)$.

Montrer que L est le point d’intersection de la droite $d$ et du plan (BGI).
	\end{enumerate}
\item On rappelle que le volume $V$ d'une pyramide est donné par la formule
\[V=\dfrac{1}{3}\times \mathcal{B}\times h\]
où $\mathcal{B}$ est l'aire d’une base et $h$ la hauteur associée à cette base.
\begin{enumerate}
\item Calculer le volume de la pyramide FBGI.
\item En déduire l'aire du triangle BGI.
\end{enumerate}
\end{enumerate}

\bigskip

\end{document}
